\section{Baseline Event Selection and Object Definitions}%
\label{ch:definitions}

In this subsection we describe the baseline event and object definitions.
The further selection optimisation is described in \cref{ch:regions}.

Electron and muons are reconstructed, identified and isolated according to the latest
recommendations by ATLAS' performance group's~\cite{web:EGammaRecommendations, web:MuonRecommendations}.
Combining electrons and muons form the leptons in this analysis. The data driven method
used by the analysis for the prediction of fake leptons in the measurement region
requires two distinct definitions of leptons, called \textit{tight}
and \textit{loose}. The method relies on measuring the transition probability
of leptons between both regions. The tight definition is the same as the final
analysis lepton definition. The transition rate from the loose to the tight
region must be small enough to ensure the stability of the methods. Also the
loose definition has to be loose enough so that the
\( N_\mathrm{tight}/N_\mathrm{loose} \) ratio is not too large and within
a small statistical uncertainty.

\subsection{Tracking and Vertex Requirements}

All of the events considered in the analysis are required to have a primary
vertex, following the guidelines of the ATLAS Inner Detector
Group~\cite{web:TrackingRecommendations}.
To further reject non-collision backgrounds originating from cosmic rays and
beam-halo events, at least one reconstructed collision vertex is required with
at least two associated tracks with transverse momentum greater than
\qty{500}{\MeV}. The vertex candidate with the highest sum of the squared
transverse momenta of all associated tracks is identified as the primary vertex.

\subsection{Trigger}%
\label{sec:definitions-triggers}

Signal trigger requirements are specific for each year of the Run 2 data
taking and are summarised in \cref{tab:trigger-analysis}. All used
triggers are unprescaled in 2015, 2016, 2017 and 2018
periods~\cite{web:LowestUnprescaled}.
The single-muon or single-electron triggers are described in
Refs.~\cite{TRIG-2016-01,TRIG-2018-05,TRIG-2018-01,ATL-DAQ-PUB-2016-001,ATL-DAQ-PUB-2017-001,ATL-DAQ-PUB-2018-002,ATL-DAQ-PUB-2019-001}.  % cspell:disable-line
A matching between online objects firing the trigger and an offline reconstructed object is required.

\begin{table}[htbp]
  \begin{center}
    \small
    \resizebox{\textwidth}{!}%
    {%
      \begin{tabular}{ l c c c }
        \toprule
                    & 2015                                                                     & 2016+2017+2018                                                                            \\
        \midrule
        \( e \)     & \texttt{HLT\_e24\_lhmedium\_L1EM20VH OR HLT\_e60\_lhmedium OR HLT\_e120\_lhloose} & \texttt{HLT\_e26\_lhtight\_nod0\_ivarloose OR HLT\_e60\_lhmedium\_nod0 OR HLT\_e140\_lhloose\_nod0} \\
        \( \mu \)   & \texttt{HLT\_mu20\_iloose\_L1MU15 OR HLT\_mu50}   & \texttt{HLT\_mu26\_ivarmedium OR HLT\_mu50}                                               \\
        \bottomrule
      \end{tabular}
    }
  \end{center}

  \caption{A summary of the trigger requirements.}%
  \label{tab:trigger-analysis}
\end{table}

All the triggers used are unprescaled High Level Triggers (HLT).
They require 1 lepton above a certain \pt threshold, which in GeV unit equals
to the number after the ``\texttt{e}'' and ``\texttt{mu}'' characters in the
trigger name. For example \texttt{HLT\_e140\_lhloose\_nod0} requires
one electron above \qty{140}{\GeV}.

Electron triggers require electrons to have at least \texttt{LHLoose},
\texttt{LHMedium} or \texttt{LHTight} offline reconstruction level, depending
on the tag in trigger name, and do not require any isolation level on the
electrons. A combination of triggers is taken where the identification criterion
is stricter going towards lower momenta to keep the trigger rate low enough.
Additionally, for the 2016, 2017 and 2018 triggers no information on the
impact parameter \dzero was used, as stated by the ``\texttt{nod0}'' tag in
the trigger name.

The low momentum muon triggers ask for an isolation requirement on the muon
according to a loose or medium working point of the muon reconstruction
algorithm at trigger level, depending on the year.

\subsection{Electrons}%
\label{sec:definitions-electrons}

The basic definition of electrons in different working points are provided
centrally by the ATLAS EGamma group~\cite{web:EGammaRecommendations} and are common to all
ATLAS analyses. These basic definitions are based on three categories:
electron reconstruction, identification, and isolation. Analyses then
choose one of the basic working points and complement it with additional
specific analysis cuts.

The reconstruction of electron objects is based on the information measured by
the EM calorimeter and the ID subsystems~\cite{ATLAS-CONF-2014-032, PERF-2013-03}.
% The electron reconstruction efficiency measured in 2015 was greater than
% \qty{98}{\%} in almost all the phase space regions~\cite{web:EGammaReco}.
In order to account for differences in simulation and data scale factors
are applied to reconstructed electrons in Monte Carlo. The scale factors
are provided centrally by the EGamma group and are expressed in bins of
electron \pt and \eta~\cite{web:EGammaRecommendations}. 
% They differ from 1.0 maximally up to
% \qty{0.5}{\%} in all of the bins~\cite{web:EGammaRecommendations}.

There is no guarantee that objects, reconstructed as electrons, truly
are electrons. The main sources of fake electron objects are jets, electrons
from photon conversions or heavy flavour hadron decays and Dalitz decays.
The EGamma group provides common working points~\cite{web:EGammaRecommendations}
to discriminate between prompt and fake electrons which are based on
a Multivariate likelihood fit of the tracking and cluster information.
The reconstruction and identification working points used are called
\texttt{LHLoose} and \texttt{LHTight}. The use of two working points
is further described in the section on the electron fakes background estimation
in \cref{sec:bkg-fakes-electron}. Briefly, the data driven method used
by the analysis to predict fake electrons in the measurement region requires two
distinct definitions of electrons (\cref{sec:bkg-fakefactor}).
They are called \textit{tight} and \textit{loose} definitions, and the methods
rely on measuring the probability of fake electrons to transition between
both regions.

The \texttt{LHTight} working point classifies a subset of electrons tagged
by the \texttt{LHLoose} working point. Similarly as for the electron
reconstruction, scale factors need to be used for
electron identification as well~\cite{web:EGammaRecommendations}.

The \texttt{PLImprovedTight} isolation~\cite{web:PLImprovedWPs} working point
is used to identify and reject the non-prompt electrons and muons candidates
from decays of heavy flavour (\(b\) and \(c\)) hadrons. It uses
a boosted decision tree algorithm trained on standard isolation variables
which measure the distance of the electron to the closest particle in a cone,
and a recurrent neural network using track impact parameters.

\begin{table}[htbp]
  \begin{center}
    \begin{tabular}{ c c c }
      \toprule
      Requirement                      & Signal electrons (tight)                  & Background electrons (loose)        \\
      \midrule
      Identification                   & \texttt{LHTight}                          & \texttt{LHLoose}                    \\
      Isolation                        & \texttt{PLImprovedTight}                  & fail \texttt{PLImprovedTight} or fail tight selection  \\
      \pt cut                          & \( \pt > \qty{10}{\GeV} \)                & \( \pt > \qty{10}{\GeV} \)           \\
      \eta cut                         & \( |\eta| < 2.47 \)                       & \( |\eta| < 2.47 \)                 \\
      Crack veto                       & veto \( 1.37 < |\eta| < 1.52 \)           & veto \( 1.37 < |\eta| < 1.52 \)     \\
      \(|\dzero|/\sigma_{\dzero}\) cut & \(|\dzero|/\sigma_{\dzero} < 5.0\)        & \(|\dzero|/\sigma_{\dzero} < 5.0\)  \\
      \(|\zzsth|\) cut                 & \(|\zzsth| < \qty{0.5}{\mm}\)             & \(|\zzsth| < \qty{0.5}{\mm}\)        \\
      Bad cluster veto                 & yes                                       & yes                                 \\
      \bottomrule
    \end{tabular}
  \end{center}

  \caption{A summary of the electron definitions in the analysis.}%
  \label{tab:definitions-electrons}
\end{table}



To avoid trigger bias and to ensure a large enough difference between tight
and loose selection of electrons in the fake estimation, the \texttt{LHTight}
identification level was chosen as the signal working point.
The \texttt{PLImprovedTight} isolation working point was chosen for the signal
electron definition. The definition of electrons in the analysis is
summarised in \cref{tab:definitions-electrons}.
The selections on the transverse \dzero and longitudinal \zzero track impact parameters are
used to ensure that the electron originates from a primary vertex.
The object quality cut ``Bad cluster veto'' is a requirement of a good pulse shape in the LAr cells
forming the cluster and is stored in the
\texttt{xAOD::EgammaParameters::BADCLUSELECTRON} flag.  % cspell:disable-line

Leptons are counted using the \qty{10}{\GeV} cut to reduce multi-lepton event
contamination of the single-lepton selection. Later in analysis regions
further cut of \qty{30}{\GeV} is applied.


\subsection{Muons}%
\label{sec:definitions-muons}

Two different sets of muon selection requirements are defined for this analysis.
Muons satisfying these criteria are referred to as \textit{tight} and
\textit{loose} muons, as defined in \cref{tab:definitions-muons}.
Tight muons are selected in the measurement regions, while loose
muons are primarily used for background estimation.

\begin{table}[htbp]
  \begin{center}
    \begin{tabular}{ c c c }
      \toprule
      Requirement                       & Signal muons (tight)                  & Background muons (loose)              \\
      \midrule
      Quality                           & \texttt{Medium}                       & \texttt{Medium}                       \\
      Bad muon veto                     & yes                                   & yes                                   \\
      Isolation                         & \texttt{PLImprovedTight}              & fail \texttt{PLImprovedTight}          \\
      \pt cut                           & \( \pt > \qty{10}{\GeV} \)            & \( \pt > \qty{10}{\GeV} \)             \\
      \eta cut                          & \( |\eta| < 2.5 \)                    & \( |\eta| < 2.5 \)                    \\
      \(|\dzero|/\sigma_{\dzero}\) cut  & \( |\dzero|/\sigma_{\dzero} < 3.0 \)  & \( |\dzero|/\sigma_{\dzero} < 3.0 \)  \\
      \(|\zzsth|\) cut                  & \( |\zzsth| < \qty{0.5}{\mm} \)        & \( |\zzsth| < \qty{0.5}{\mm} \)        \\
      \bottomrule
    \end{tabular}
  \end{center}

  \caption{Summary of the muon definitions in the analysis.}%
  \label{tab:definitions-muons}
\end{table}

Muons with \(\pt > \qty{10}{\GeV}\) and within the acceptance \(|\eta| < 2.5\)
are used in the analysis. Quality requirements are imposed on the muon
reconstruction and identification according to the \texttt{Medium} working point~\cite{PERF-2015-10}.

Background muons are real muons embedded within a jet and produced by in-flight
decays of secondary mesons inside the jet. For this reason, the detector
activity surrounding these muons, known as \textit{isolation}, is used
to discriminate between background and prompt muons.
Signal muons must fulfil the criteria of the \texttt{PLImprovedTight} isolation working 
point~\cite{web:PLImprovedWPs},
while loose muons are required to fail this selection.

Finally, requirements on the transverse impact parameter significance,
\( \mathrm{Sig}(\dzero)=|\dzero|/\sigma_{\dzero} \),
and the longitudinal impact parameter multiplied by the polar angle,
\( |\zzsth| \), are utilised to distinguish between prompt and background muons,
as described in \cref{tab:definitions-muons}.

To account for isolation and reconstruction selection efficiency differences
between data and MC, the muon performance group of the ATLAS collaboration
provides calibration factors measured in data which are applied to simulated
muons in this analysis~\cite{web:MuonRecommendations}.

Leptons are counted using the \qty{10}{\GeV} cut to reduce multi-lepton event
contamination of the single-lepton selection. Later in analysis regions
further cut of \qty{30}{\GeV} is applied.


\subsection{Jets}%
\label{sec:definitions-jets}

Jets or particles originating from the hadronisation of partons are
reconstructed by clustering energy deposits in the calorimeter calibrated at
the EM scale. The \antikt algorithm~\cite{Cacciari:2008gp} is used with
a radius parameter of 0.4. The signal selection requires 
\(\pt > \qty{25}{GeV}\) jets with no $\eta$ requirements but the subsequent \bjet
requirement restricts our jets to be central with \(\abseta < 2.5\),
as described below.
An additional category of jets, the baseline jets are defined with more loose
selection criteria. The baseline jets are used for the \met reconstruction (\cref{sec:definitions-met}) and
overlap removal procedure (\cref{sec:definitions-overlap-removal}).
For all jets the expected average transverse energy
contribution from pile-up is subtracted using an area-based \pt density
subtraction method and a residual correction derived from the MC simulation,
both detailed in \Refn{\cite{Cacciari:2007fd, PERF-2016-04}}.
The jet definitions and selections are summarised in \cref{tab:definitions-jets}.

\begin{table}[htbp]
  \begin{center}
    \begin{tabular}{ c c c }
      \toprule
      Requirement         & Signal jets                                   & Baseline jets                \\
      \midrule
      Jet type            & \texttt{AntiKt4EMPFlowJets}                   & \texttt{AntiKt4EMPFlowJets}  \\
      JVT working point   & \texttt{Tight}                                & \texttt{Tight}               \\
      fJVT working point  & ---                                           & ---                          \\
      \pt cut             & \( \pt > \qty{25}{\GeV} \)                    & \( \pt > \qty{20}{\GeV} \)    \\
      \eta cut            & \( |\eta| < 2.5 \)                            & \( |\eta| < 4.5 \)           \\
      \(b\)-tagging       & \texttt{DL1r} with \texttt{FixedCutBEff\_77}  &                              \\
      \bottomrule
    \end{tabular}
  \end{center}

  \caption{Summary of the jet definitions in the analysis.}%
  \label{tab:definitions-jets}
\end{table}

The majority of pile-up jets are rejected using the \textit{jet vertex tagger}
(JVT) tool~\cite{ATLAS-CONF-2014-018}, which is a combination of track-based
variables providing discrimination against pile-up jets. The \texttt{Tight}
working point is used for jets in the momentum range \( 20 < \pt < \qty{60}{\GeV} \)
and \( |\eta| < 2.4 \).

In this analysis, events containing jets tagged as \bjets are selected,
as described in \cref{ch:regions}. These jets are identified as
detailed in~\cite{ATL-PHYS-PUB-2016-012}. Due to the limitations of the identification
of \bjets only jets within \(|\eta| < 2.5\) are considered.
The \texttt{DL1r} tagger~\cite{web:FlavourTagging} was used with a fixed-cut
working point selecting \bjets with \qty{77}{\%} efficiency.

Additionally, the jet-based event cleaning procedure~\cite{web:JetCleaning},
recommended by the Jet/EtMiss group was performed. The \textit{loose} working
point was chosen.


\subsection{Missing Transverse Momentum}%
\label{sec:definitions-met}

The missing transverse momentum, or \textit{missing transverse energy} \MET,
arises due to undetected energy in the event, either due to neutrinos escaping
detection or to other particles outside the acceptance, badly reconstructed
or failing reconstruction.

The \MET is computed as the unbalance of the total visible \pt in the event with
contributions from all visible objects in the event which are selected,
reconstructed and fully calibrated, plus a term, the so-called
\textit{soft-term}, accounting for all tracks in the event which are not
associated with any visible object but are associated with the identified
hard-scatter vertex. The usage of a track-based soft-term is motivated by an
improved performance in \MET reconstruction against pile-up~\cite{JETM-2020-03}.
The \texttt{Tight} working point is used for the \MET calculation.


\subsection{Overlap Removal}%
\label{sec:definitions-overlap-removal}

% description of the procedure (standard)
% https://indico.cern.ch/event/539619/contributions/2191033/attachments/1284239/1909241/Farrell_ort_asg.pdf

It could happen that a physical object is reconstructed by several different
algorithms designed to reconstruct different particles, e.g.\ an electron is
also reconstructed as a jet, resulting in a double counting of energy deposit in
the calorimeter. To account for this, ambiguous objects are removed following
the recommendations of the study documented in reference~\cite{Adams:1700874}
and by using the produced common tool~\cite{web:OverlapRemoval}.
The `Standard' working-point of the tool is used and electrons, muons, and jets
are considered by the algorithms in all regions.
Loose object definitions are used for the procedure.

The overlap removal algorithm resolves overlaps in the following order:
\begin{itemize}
  \item \textbf{muon-PFlow jets:} Removes PFlow jets that are actually muons.
  \item \textbf{electron-muon:} If a muon overlapping with an electron leaves
  a sufficiently high energy deposit in the calorimeter and shares a track
  reconstructed in the inner detector with the electron, then the muon is removed.
  Any electron sharing an ID track with the remaining muons is removed.
  \item \textbf{electron-jet:} Jets within \(\Delta R < 0.2\) of an electron
  are removed. Electrons within \(\Delta R < 0.4\) of the remaining jets are
  discarded.
  \item \textbf{muon-jet:} Jets within \(\Delta R < 0.2\) of a muon
  and either featuring less than three tracks or fulfilling
  \(\pt(\mu) / \pt(\mathrm{jet}) > 0.5\) and
  \(\pt(\mu) / \sum\pt(\mathrm{all\;jet\;tracks}) > 0.7\) are removed.
  Muons overlapping with \(\Delta R < 0.4\) of the remaining jets are removed.
\end{itemize}

\subsection{Particle level object definition}

Stable particles are used for the definition of particle level objects.
Every particles with a lifetime \( \tau \) such that \(c\tau > \qty{10}{mm}\)
is fulfilled, is considered a stable particle.
The objects are defined in a way to be as close to the detector level as possible.

Leptons are defined as electrons, muons and neutrinos originating from \Wboson
or \Zboson decays, as well as from \( \tau \)-lepton decays.
The charged leptons are ``dressed'' by recombining them with every photon
within a \(\Delta R = 0.1\) cone.
They are counted using the \qty{10}{\GeV} and the pseudo-rapidity
\(|\eta| < 2.5\) cuts to reduce multi-lepton event
contamination of the single-lepton selection.
Additionally, they have to pass cuts on the
transverse momentum \(\pt > \qty{30}{\GeV}\) in the analysis regions.

The particle level small-\(R\) jet collection is defined using the four-momenta
of stable particles. The leptons defined above are excluded from jets.
Non-prompt leptons from hadron decays are included in the clustering.
We use the same \antikt algorithm with the radius parameter \(0.4\) as the
detector level jets defined above. In the ATLAS software, this jet collection
is denoted as \texttt{AntiKt4TruthDressedWZJets}.
Only jets with \(\pt > \qty{25}{\GeV}\) and \(|\eta| < 2.5\) are considered
for the analysis. We do not apply any overlap removal on particle level.

A jet is \btagged if a \(b\)-hadron with \(\pt > \qty{5}{\GeV}\) can
be found within \(\Delta R (jet, b) < 0.3\). The hadron is assigned to
the closest jet, in case two jets satisfy the \(\Delta R\) requirement.
This corresponds to the \texttt{HadronConeExclTruthLabelID} variable in
the ATLAS software.

Particle-level selections follow closely the one of the analysis. This also implies that events with two leptons 
(e.g.\ dilepton \ttbar events) may be selected in the fiducial volume if one lepton falls outside the detector acceptance. 
The composition of reconstructed events based on true \(WW\) decay products
is shown in \cref{fig:WWcomposition}.

\begin{figure}[htbp]
  \centering
  \subfloat[]{
    \includegraphics[width=0.3\textwidth]{Composition/IR}%
    \label{fig:WWcomposition:IR}
  }
  \subfloat[]{
    \includegraphics[width=0.3\textwidth]{Composition/SR}%
    \label{fig:WWcomposition:SR}
  }
  \subfloat[]{
    \includegraphics[width=0.3\textwidth]{Composition/WR}%
    \label{fig:WWcomposition:WR}
  }

  \caption{True \(WW\) composition in the reconstructed events for
    \protect\subref{fig:WWcomposition:IR} inclusive bulk selection,
    \protect\subref{fig:WWcomposition:SR} search-like selection, and
    \protect\subref{fig:WWcomposition:WR} hadronic \Wboson selection.
    All events are shown including background. An event is classified as `Unknown'
    if the number of \Wboson bosons is different than two.
  }%
  \label{fig:WW:composition}
\end{figure}

\subsection{Software Infrastructure}

\texttt{AnalysisBase} version 21.2.242 is used to process the data and calibrate
it. CP algorithms\footnote{\url{https://gitlab.cern.ch/tadej/MultiLeptonAnalysis}}
are used for the ntuples production and have been validated against
\texttt{AnalysisTop}.\ \texttt{TNAnalysis}\footnote{\url{https://gitlab.cern.ch/tadej/TNAnalysis}}
is used for final processing and plotting.
